\scanpage{225}
\mspage{175}
% L1–6
\deleted{Comme $X_{/x}\to e$ est dans $W$, il s'ensuit que $\Omega_{/x}\to\widetilde\Omega(x)$ l'est aussi (moyennant W\unsure{17|7|1} bis). D'autre part, c'est là une flèche \deleted{des} d'\deleted{catégories}\inserted{objets} \uline{ponctués} de $\mathrm{Cat}$, $\Omega_{/x}$ étant ponctué comme on vient de le dire, $\widetilde\Omega(x)$ \unsure{par} $e\longrightarrow\widetilde\Omega(x)$ du diagramme$^*$.}
% L6 left margin
% slanted
\margin{Non, il vaut mieux le ponctuer par l'image du pt marqué de $\deleted{X}\Omega_{/x}$ \textendash{} \unsure{Attention} que le carré n'est pas commutatif.}
% L6–7
\deleted{En résumé, on a des flèches dans $\mathrm{Cat}_{\bullet}$}
% L8
\deleted{\[
\Omega_x\longrightarrow\Omega_{/x}\longrightarrow\widetilde\Omega(x)
\]}
% L9–14
\deleted{qui sont toutes deux dans $W$ en tant que flèches de $\mathrm{Cat}$, donc définissent des iso. dans $\mathrm{Hot}_{\bullet W}$. Ni le premier ni le deuxième terme n'est fonctoriel en $x$ (i.e. ne \uline{définit un foncteur $X\longrightarrow\mathrm{Cat}_\bullet$), mais}}
% L14–17
de l'image $\sigma(\unsure{x})_{/x'}$ du pt marqué de $\Omega_{/x}$ vers le point marqué $\sigma(x')_{/x'}$ de $\Omega_{/x'}$. Ceci doit suffire pour définir la flèche
% L18
\[
(\Omega_{/x},\sigma(\unsure{x})_{/x})\unsure{-|\longrightarrow}(\Omega_{/x'},\sigma(x')_{/x'})\quad\deleted{\text{\unsure{de}}}
\]
% L19–21
comme flèche dans $\mathrm{Hot}_{\bullet W}$, avec transitivité pour une composition $x\xrightarrow{\unsure{u}}x'\xrightarrow{v}x''$ dans $X$. Je vais pour \unsure{l'instant} me borner
% archive sheet number: 225
